Er zijn verschillende manieren om te achterhalen welke netwerk devices we hebben. We geven er hier twee:
\begin{lstlisting}[language=bash]
cat /proc/net/dev
\end{lstlisting}

\begin{lstlisting}[language=bash]
ip link show
\end{lstlisting}

In de output van beide commando's zien we de interface \texttt{lo} wat de localhost interface is en we hebben een interface die begint met \texttt{en} wat een ethernet interface is.

Na een
\begin{lstlisting}[language=bash]
cat /etc/network/interfaces
\end{lstlisting}
zien we een output die lijkt op:
\begin{lstlisting}[language=bash]
# This file describes the network interfaces available on your system
# and how to activate them. For more information, see interfaces(5).

source /etc/network/interfaces.d/*

# The loopback network interface
auto lo
iface lo inet loopback

# The primary network interface
allow-hotplug ens33
iface ens33 inet dhcp
\end{lstlisting}
We zien dat we in de twee regel van het bestand verwezen worden naar de manual-pagina van interfaces. Daar staat veel meer in. Lees deze man-page door tot en met het stukje over INCLUDING OTHER FILES.

Verder zien we in ons \texttt{interfaces} bestand nog twee regels die beginnen met \texttt{iface}. Na de \texttt{iface} optie komt eerst de naam van de interface die we willen configureren. Het key-word \texttt{inet} geeft aan dat we het hebben over een Internet-configuratie. Het Internet gebruikt TCP/IP en dat is wat hier bedoeld wordt, inet = IP. Tot slot zien we bij de \texttt{lo} interface dat er \texttt{loopback} gebruikt wordt. De \texttt{loopback} methode kent geen verdere opties en zorgt ervoor dat op de \texttt{lo} interface het 127.0.0.1 IP adres komt te staan. Voor de \texttt{ens33} interface, die bij jou best anders kan heten, zien we de methode \texttt{dhcp}. Op deze interface zal het systeem dus proberen via DHCP zijn settings te verkrijgen.

In de man-page van interfaces hebben we ook gezien dat we een methode \texttt{static} hebben waarmee een statisch IP-adres op de interface kunnen zetten.
\begin{lstlisting}[language=bash]
iface eth1 inet static
     address 192.168.1.2/24
     gateway 192.168.1.1
\end{lstlisting}
Static kent een aantal opties die we meekunnen geve. Let op het inspringen! De ingesprongen opties behoren bij de configuratie van \texttt{eth1}. De \texttt{eth1} interface krijgt het IP-adres 192.168.1.2 met als subnetmask 255.255.255.0. Ook wordt de default gateway gezet naar 192.168.1.1.

